% test-009.tex -- comandos \spdiv y \spmul
%
% Compilar:  lualatex-dev test-009.tex
% Validar:   show-pdf-tags --xml test-009.pdf | rnv-wrapp
%            verapdf --flavour ua2 --format text test-009.pdf
%
% Cada caso es un parrafo numerado, para poder referirse a el al escuchar
% el PDF. Los comentarios "doc:" indican lo que documenta el manual; en el
% resto, la lectura se revisa escuchando. Ambos comandos trabajan en modo
% matematico y pasan sus dos argumentos por \spnum, asi que solo reciben
% numeros. Las entradas invalidas (argumentos vacios o que no son numeros,
% que detienen la compilacion) no van aqui.
\DocumentMetadata
  {
    lang=es-CL,
    pdfversion=2.0,
    pdfstandard=ua-2,
    tagging=on,
    tagging-setup={math/setup={mathml-SE}},
  }
\documentclass{article}
\usepackage[spanish]{babel}
\usepackage{unicode-math}
\usepackage{spintent}
\usepackage{hyperref}
\hypersetup
  {
    pdfdisplaydoctitle = true,
    pdftitle           = {test-009: spdiv y spmul},
  }
\pagestyle{empty}
\begin{document}

\section*{A. Comando spdiv: argumentos}

% doc: «20 dividido 26»
1. Dos enteros: $\spdiv{20}{26}$.

% doc: «menos 20 dividido 26»
2. Primer argumento negativo: $\spdiv{-20}{26}$.

% doc: «20 dividido menos 26»
3. Segundo argumento negativo: $\spdiv{20}{-26}$.

% doc: «menos 20 dividido menos 26»
4. Ambos negativos: $\spdiv{-20}{-26}$.

% doc: «más 20 dividido 26»
5. Con signo positivo explícito: $\spdiv{+20}{26}$.

% doc: «3 coma 5 dividido 2»
6. Decimales con coma: $\spdiv{3,5}{2}$.

% doc: «3 punto 5 dividido 2»
7. Decimales con punto: $\spdiv{3.5}{2}$.

% doc: «1 coma 3 periódico dividido 2»
8. Número periódico: $\spdiv{1;3}{2}$.

% doc: «3 coma 1 con 4 periódico dividido 2»
9. Periódico mixto: $\spdiv{3,1;4}{2}$.

% doc: «3 coma 14 por 10 al cubo dividido 2»
10. Notación científica: $\spdiv{3,14e3}{2}$.

% doc: «100000 dividido 250»
11. Cifras grandes: $\spdiv{100000}{250}$.

% doc: «1234567 coma 891234 dividido 10»
12. Decimales grandes: $\spdiv{1234567,891234}{10}$.

% doc: «2 dividido 26»
13. Primero menor que el segundo: $\spdiv{2}{26}$.

% doc: «0 dividido 5»
14. Con cero: $\spdiv{0}{5}$.

\section*{A. Comando spdiv: llaves de spnum}

% doc: «100000 dividido 250»
15. Con cifras-sep false: $\spdiv[cifras-sep=false]{100000}{250}$.

% doc: «100000 dividido 250»
16. Con cifras-sep false y espacios del usuario: $\spdiv[cifras-sep=false]{10\,00\,00}{250}$.

% doc: «20 negativo dividido 26»
17. Con read-sign clear en el primero: $\spdiv[read-sign=clear]{-20}{26}$.

% doc: «20 dividido 26 negativo»
18. Con read-sign clear en el segundo: $\spdiv[read-sign=clear]{20}{-26}$.

% doc: «1 punto 3 periódico dividido 2»
19. Con read-period-sep punto: $\spdiv[read-period-sep=punto]{1;3}{2}$.

% doc: «3 coma 14 por 10 al cubo dividido 2»
20. Con notation-sym times: $\spdiv[notation-sym=times]{3,14e3}{2}$.

% doc: «3 coma 14 negativo por 10 al cubo dividido 2»
21. Llaves de spnum combinadas: $\spdiv[read-sign=clear,notation-sym=times]{-3,14e3}{2}$.

\section*{A. Comando spdiv: símbolo y lectura}

% doc: «20 dividido 26»
22. Con set-sym old-style: $\spdiv[set-sym=old-style]{20}{26}$.

% doc: «20 dividido 26»
23. Con set-sym two-dots: $\spdiv[set-sym=two-dots]{20}{26}$.

% doc: «20 dividido 26»
24. Con set-sym slash: $\spdiv[set-sym=slash]{20}{26}$.

% doc: «20 dividido 26»
25. Con read-sym dividido: $\spdiv[read-sym=dividido]{20}{26}$.

% doc: «20 dividido por 26»
26. Con read-sym dividido-por: $\spdiv[read-sym=dividido-por]{20}{26}$.

% doc: «20 dividido entre 26»
27. Con read-sym dividido-entre: $\spdiv[read-sym=dividido-entre]{20}{26}$.

% doc: «20 dividido por 26»
28. Combinando old-style y dividido-por: $\spdiv[set-sym=old-style,read-sym=dividido-por]{20}{26}$.

% doc: «20 dividido entre 26»
29. Combinando slash y dividido-entre: $\spdiv[set-sym=slash,read-sym=dividido-entre]{20}{26}$.

\section*{A. Comando spdiv: paréntesis}

% doc: «menos 20 dividido 26»
30. Con wrap-parens false, el valor por defecto: $\spdiv[wrap-parens=false]{-20}{26}$.

% doc: «se abren paréntesis menos 20 dividido 26 se cierran paréntesis»
31. Con wrap-parens true: $\spdiv[wrap-parens=true]{-20}{26}$.

% doc: «se abren paréntesis menos 20 dividido 26 se cierran paréntesis»
32. Con wrap-parens true y read-parens true: $\spdiv[wrap-parens=true,read-parens=true]{-20}{26}$.

% doc: «menos 20 dividido 26»
33. Con wrap-parens true y read-parens false: $\spdiv[wrap-parens=true,read-parens=false]{-20}{26}$.

% doc: «menos 20 dividido 26»
34. Con wrap-parens false y read-parens false: $\spdiv[wrap-parens=false,read-parens=false]{-20}{26}$.

% doc: «se abren paréntesis menos 20 dividido 26 se cierran paréntesis»
35. Con paréntesis y slash: $\spdiv[wrap-parens=true,set-sym=slash]{-20}{26}$.

% doc: «20 negativo dividido por 26»
36. Todas las llaves juntas: $\spdiv[wrap-parens=true,read-parens=false,set-sym=old-style,read-sym=dividido-por,read-sign=clear]{-20}{26}$.

\section*{B. Comando spmul: argumentos}

% doc: «20 por 26»
37. Dos enteros: $\spmul{20}{26}$.

% doc: «menos 20 por 26»
38. Primer argumento negativo: $\spmul{-20}{26}$.

% doc: «20 por menos 26»
39. Segundo argumento negativo: $\spmul{20}{-26}$.

% doc: «menos 20 por menos 26»
40. Ambos negativos: $\spmul{-20}{-26}$.

% doc: «más 20 por 26»
41. Con signo positivo explícito: $\spmul{+20}{26}$.

% doc: «3 coma 5 por 2»
42. Decimales con coma: $\spmul{3,5}{2}$.

% doc: «3 punto 5 por 2»
43. Decimales con punto: $\spmul{3.5}{2}$.

% doc: «1 coma 3 periódico por 2»
44. Número periódico: $\spmul{1;3}{2}$.

% doc: «3 coma 1 con 4 periódico por 2»
45. Periódico mixto: $\spmul{3,1;4}{2}$.

% doc: «3 coma 14 por 10 al cubo por 2»
46. Notación científica: $\spmul{3,14e3}{2}$.

% doc: «100000 por 250»
47. Cifras grandes: $\spmul{100000}{250}$.

% doc: «1234567 coma 891234 por 10»
48. Decimales grandes: $\spmul{1234567,891234}{10}$.

% doc: «2 por 26»
49. Primero menor que el segundo: $\spmul{2}{26}$.

% doc: «0 por 5»
50. Con cero: $\spmul{0}{5}$.

\section*{B. Comando spmul: llaves de spnum}

% doc: «100000 por 250»
51. Con cifras-sep false: $\spmul[cifras-sep=false]{100000}{250}$.

% doc: «100000 por 250»
52. Con cifras-sep false y espacios del usuario: $\spmul[cifras-sep=false]{10\,00\,00}{250}$.

% doc: «20 negativo por 26»
53. Con read-sign clear en el primero: $\spmul[read-sign=clear]{-20}{26}$.

% doc: «20 por 26 negativo»
54. Con read-sign clear en el segundo: $\spmul[read-sign=clear]{20}{-26}$.

% doc: «1 punto 3 periódico por 2»
55. Con read-period-sep punto: $\spmul[read-period-sep=punto]{1;3}{2}$.

% doc: «3 coma 14 por 10 al cubo por 2»
56. Con notation-sym times: $\spmul[notation-sym=times]{3,14e3}{2}$.

% doc: «3 coma 14 negativo por 10 al cubo por 2»
57. Llaves de spnum combinadas: $\spmul[read-sign=clear,notation-sym=times]{-3,14e3}{2}$.

\section*{B. Comando spmul: símbolo y lectura}

% doc: «20 por 26»
58. Con set-sym cross: $\spmul[set-sym=cross]{20}{26}$.

% doc: «20 por 26»
59. Con set-sym dot: $\spmul[set-sym=dot]{20}{26}$.

% doc: «20 por 26»
60. Con set-sym asterisk: $\spmul[set-sym=asterisk]{20}{26}$.

% doc: «20 por 26»
61. Con read-sym por: $\spmul[read-sym=por]{20}{26}$.

% doc: «20 multiplicado por 26»
62. Con read-sym multiplicado-por: $\spmul[read-sym=multiplicado-por]{20}{26}$.

% doc: «20 veces 26»
63. Con read-sym veces: $\spmul[read-sym=veces]{20}{26}$.

% doc: «20 multiplicado por 26»
64. Combinando cross y multiplicado-por: $\spmul[set-sym=cross,read-sym=multiplicado-por]{20}{26}$.

% doc: «20 veces 26»
65. Combinando asterisk y veces: $\spmul[set-sym=asterisk,read-sym=veces]{20}{26}$.

\section*{B. Comando spmul: paréntesis}

% doc: «menos 20 por 26»
66. Con wrap-parens false, el valor por defecto: $\spmul[wrap-parens=false]{-20}{26}$.

% doc: «se abren paréntesis menos 20 por 26 se cierran paréntesis»
67. Con wrap-parens true: $\spmul[wrap-parens=true]{-20}{26}$.

% doc: «se abren paréntesis menos 20 por 26 se cierran paréntesis»
68. Con wrap-parens true y read-parens true: $\spmul[wrap-parens=true,read-parens=true]{-20}{26}$.

% doc: «menos 20 por 26»
69. Con wrap-parens true y read-parens false: $\spmul[wrap-parens=true,read-parens=false]{-20}{26}$.

% doc: «menos 20 por 26»
70. Con wrap-parens false y read-parens false: $\spmul[wrap-parens=false,read-parens=false]{-20}{26}$.

% doc: «se abren paréntesis menos 20 por 26 se cierran paréntesis»
71. Con paréntesis y asterisk: $\spmul[wrap-parens=true,set-sym=asterisk]{-20}{26}$.

% doc: «20 negativo multiplicado por 26»
72. Todas las llaves juntas: $\spmul[wrap-parens=true,read-parens=false,set-sym=cross,read-sym=multiplicado-por,read-sign=clear]{-20}{26}$.

\section*{C. En fórmulas}

% doc: «12 dividido 4 es igual a 3»
73. División con igualdad: $\spdiv{12}{4} = 3$.

% doc: «3 por 4 es igual a 12»
74. Multiplicación con igualdad: $\spmul{3}{4} = 12$.

% doc: «12 dividido 4 más 3 por 4»
75. Suma de una división y un producto: $\spdiv{12}{4} + \spmul{3}{4}$.

% doc: «se abren paréntesis menos 12 dividido 4 se cierran paréntesis es igual a menos 3»
76. División con paréntesis: $\spdiv[wrap-parens=true]{-12}{4} = -3$.

% doc: «12 dividido 4 es igual a 12 dividido 4»
77. Llaves distintas en una misma fórmula: $\spdiv[set-sym=slash]{12}{4} = \spdiv[set-sym=old-style]{12}{4}$.

% doc: «3 por 4 dividido 2 es igual a 6»
78. Producto y división juntos: $\spmul{3}{4} \spdsym 2 = 6$.

% doc: «se abren paréntesis menos 3 por 4 se cierran paréntesis es igual a menos 12»
79. En fórmula destacada:
\[ \spmul[wrap-parens=true]{-3}{4} = -12 \].

\end{document}
